\documentclass[track=junior,level=1]{kidscode}
\begin{document}
\kctitle{Notes for Families}{What we are doing, how we keep children safe, and how you can help at home}

\section{Dear families,}
Welcome to Little Robots! Once a week, for 45 minutes, your child will learn the first ideas of
coding, in a small group of up to four children (about 7--9 years old, US Grades 2--3). No
experience is needed, and your child does not need to read well yet: we read every page aloud.

In ten weeks children learn that a computer follows exact steps, in order. They start by being
robots themselves on a grid of tape on the floor. Then they use the free \textbf{Scratch} app to
snap colored blocks together and make a cat talk, dance, react to a click and draw with a pen.
\begin{itemize}[itemsep=2pt]
  \item \textbf{Weeks 1--4:} robots and steps, meeting Scratch, fixing bugs, repeat loops.
  \item \textbf{Week 5:} a play day.
  \item \textbf{Weeks 6--8:} clicks and keys, two sprites talking, drawing with the pen.
  \item \textbf{Weeks 9--10:} the project, \enquote{My Animated Greeting Card}, and a showcase
        where your child shows the card to the family and gets a certificate.
\end{itemize}
Every session has the same five parts: a warm-up game or short video, a short demonstration
(\enquote{Show me}), building in pairs (\enquote{Make it}), a wiggle break, and a Share where each
child shows what they made. Children work as a pair: the \textbf{Driver} holds the mouse and the
\textbf{Helper} points and reads. They swap when the sand timer runs out.

\section{Our safety and privacy promises}
\begin{safety}[What we promise]
\begin{itemize}[itemsep=2pt]
  \item \textbf{Offline.} We use the Scratch app on the computer, not the website. Nothing needs the internet.
  \item \textbf{No accounts.} Nobody signs up, logs in or gives an email address.
  \item \textbf{Nicknames only.} Projects show a nickname or a first name, never a full name,
        address, school, age or photo.
  \item \textbf{Nothing is uploaded.} Projects stay on the computer. The showcase is at home, with
        family only, on the screen.
  \item \textbf{Ask a grown-up.} Children learn to stop and ask whenever an app asks for anything.
  \item \textbf{Healthy screen time.} Every session has time away from the screen and a wiggle break.
  \item \textbf{For everyone.} All pages use color-blind-safe colors with words and icons, and every
        page comes in a dyslexia-friendly edition. Please ask.
\end{itemize}
\end{safety}

\section{What your child needs}
\begin{itemize}[itemsep=2pt]
  \item \textbf{A Mac with Scratch.} The Scratch 3 app is free in the Mac App Store. No account is needed.
  \item \textbf{A mouse.} Dragging blocks with a trackpad is hard for small hands.
  \item \textbf{45 minutes, once a week,} at a time when your child is not too tired or hungry.
  \item \textbf{A grown-up nearby,} to read the pages aloud and to be the Helper when there is one child.
\end{itemize}

\section{How to help at home}
You do not need to know anything about coding to help.
\begin{itemize}[itemsep=3pt]
  \item \textbf{Play the play-at-home card.} Each week there is an optional 5--10 minute game, mostly
        without a computer. Stop while it is still fun.
  \item \textbf{Read the pages aloud.} Point at the pictures and let your child point too. Saying,
        drawing and circling are all good answers.
  \item \textbf{Praise effort.} \enquote{You kept trying until it worked!} means more than
        \enquote{You're so clever.}
  \item \textbf{Bugs are normal.} A mistake in code is called a bug, and finding one is a real skill.
        Ask \enquote{What did you want it to do? What did it do instead?}
  \item \textbf{Hands off the mouse.} If you help, be the Helper: point and use words, and let your child click.
\end{itemize}

\section{Ask your child about\dots}
One question for each week. They make great conversation starters at dinner.

\begin{center}\small\renewcommand{\arraystretch}{1.45}
\begin{tabularx}{\linewidth}{@{}>{\bfseries}p{1.0cm}>{\RaggedRight}p{4.2cm}>{\RaggedRight\arraybackslash}X@{}}
\toprule
Week & Topic & Ask your child\dots\\
\midrule
1 & Robots Need Instructions & \enquote{How did you make the robot walk to the goal? What happened when a card was in the wrong place?}\\
2 & Meet Scratch & \enquote{What is a sprite? What did your cat say?}\\
3 & Bugs Are Clues & \enquote{What is a bug? How did you find one today?}\\
4 & Repeat It! & \enquote{Can you show me a dance with a repeat in it?}\\
5 & Play Day and Show Me & \enquote{What are you best at in Scratch now?}\\
6 & When I Click\dots & \enquote{What happens when you click your sprite or press a key?}\\
7 & Two Sprites Talk & \enquote{How did you make two sprites take turns to talk?}\\
8 & Draw With the Pen & \enquote{How does the cat draw a square? Can you walk one for me?}\\
9 & Project Studio & \enquote{Who is your greeting card for? What happens at the start, middle and end?}\\
10 & Showcase and Certificate & \enquote{What are you proudest of? What would you like to make next?}\\
\bottomrule
\end{tabularx}
\end{center}

\section{How to celebrate}
\begin{itemize}[itemsep=3pt]
  \item \textbf{Come to the showcase} in week 10. Your child will show the greeting card, tell you what
        one script does and maybe show you a bug they fixed. You get sticker feedback cards to fill in.
  \item \textbf{Certificates.} Every child who finishes Little Robots gets a printed certificate.
  \item \textbf{A fridge gallery.} Ask your child to draw their card or their favorite sprite for the fridge.
        Any photos or screenshots stay at home, not online.
  \item \textbf{Celebrate effort.} \enquote{You found that bug all by yourself!} is the best thing a
        young coder can hear.
\end{itemize}

\bigskip
Thank you for coding with us. If anything would help your child learn (for example bigger text,
more breaks or a different time of day), please tell me.

\bigskip
With warm wishes,\par\bigskip
\rule{6cm}{0.4pt}\par
{\small (Your coding club host)}
\end{document}
